\chapter{Leakage: The Boss Fight}
\label{ch:leakage}

At $R^2 \approx 0.01$, a leak that adds $0.001$ of fake signal --- one part
in a thousand of the target's variance --- inflates your validation score by
ten percent and evaporates on the private leaderboard. Leakage at low SNR is
not one risk among many; it is the dominant failure mode, the first
explanation you should reach for whenever a result looks good, and the
reason this book devotes a full chapter to what other treatments cover in a
paragraph.

\section{What leakage is, precisely}

A model leaks when its inputs at prediction time $t$ contain information
that would not be available at time $t$ under the deployment protocol.
The definition has two halves that fail independently:

\begin{itemize}
\item \textbf{Temporal leakage}: an input encodes data from times $> t$.
  Classic vectors: centered rolling windows, forward-filled-then-shifted
  columns, target encodings computed over the full dataset, scalers fit on
  all rows.
\item \textbf{Protocol leakage}: an input encodes data from times $\le t$
  that the \emph{deployment protocol} nevertheless withholds. This is the
  subtle half. In this competition you never observe today's responders,
  at any earlier timestamp of today --- even though they are ``past''
  data. A within-day lag of the target is temporally causal and still
  illegal.
\end{itemize}

The second half is why the first commandment of this genre is:

\begin{keybox}[: model the deployment information set explicitly]
Write down, as a formal statement, exactly what is knowable at prediction
time. For this competition: \emph{at (day $d$, time $t$) you know: all
features of all symbols for day $d$ up to and including time $t$; all
features and all responders for all days $< d$; nothing else.} Every
feature must come with a one-line proof of membership in this set. If you
cannot write the proof, the feature does not ship.
\end{keybox}

\section{A type system for features}

We found it useful to force every engineered column into one of three
types, checked at code-review time like a type annotation:

\begin{center}
\begin{tabularx}{\textwidth}{@{}lXX@{}}
\toprule
type & definition & examples from the campaign \\
\midrule
\textbf{trailing} & function of data strictly before $t$ (possibly
crossing days) & per-symbol rolling mean/std over the last 1000 steps;
yesterday's responder path (the lag join); rsig momentum differences \\
\textbf{contemporaneous} & function of data delivered \emph{at} $t$ &
cross-symbol market average of a feature at (day, time) --- legal here
because the gateway delivers all symbols' rows for a timestamp together \\
\textbf{forward} & function of data after $t$ & synthetic auxiliary
targets like $\resp{6}_{t} + \resp{6}_{t+20} + \resp{6}_{t+40}$ ---
\emph{legal only as training targets, never as inputs} \\
\bottomrule
\end{tabularx}
\end{center}

The discipline pays in review speed: a pull request adding features is
approved by checking each column's declared type against its
implementation, not by re-deriving the whole pipeline's causality.

Two implementation patterns are worth stealing directly:

\begin{itemize}
\item \textbf{Lags as self-joins, not shifts.} To give day $d$ access to
  day $d{-}1$'s responders, do not \code{shift()} a sorted frame ---
  sort-order bugs and group-boundary bleed are silent. Instead, take the
  responder table, \emph{add one to its date column}, and left-join on
  (symbol, date, time). The join is order-independent, missing previous
  days produce explicit nulls, and the construction is \emph{structurally}
  incapable of reaching forward.
\begin{lstlisting}
lagged = (df.select([SYM, DATE, TIME, *resp_cols])
            .with_columns((pl.col(DATE) + 1).alias(DATE))
            .rename({c: f"{c}_lag1d" for c in resp_cols}))
df = df.join(lagged, on=[SYM, DATE, TIME], how="left")
\end{lstlisting}
\item \textbf{Trailing windows by construction.} Every rolling operation in
  the codebase goes through one helper that only exposes trailing
  semantics. Nobody types \code{center=True} because the parameter does not
  exist in the internal API. (We adopted this after the near-miss described
  below.)
\end{itemize}

\section{Preprocessing is part of the model}

The most common leak in public notebooks is not a feature at all --- it is
the scaler. Standardization statistics (means, variances, clip quantiles)
computed over \emph{all} rows transport information from the validation
period into training-time transformations. The effect is small per column
and large in aggregate, and it flatters exactly the models you are trying
to compare.

The rule: \textbf{fit all preprocessing on training dates only; transform
validation with frozen statistics; at streaming inference, update running
statistics using only past days.} Our pipeline freezes a quantile clipper
(fitted $[0.1\%, 99.9\%]$ bounds) and a standardizer at the training
boundary; during walk-forward evaluation the standardizer is allowed
\code{partial\_fit} on each day \emph{after} it has been predicted ---
mirroring exactly what a deployed system could do.

The same boundary governs \emph{data-dependent feature definitions}: the
16-feature subset selected for rolling statistics, clustering assignments,
symbolic-combination choices (Chapter~\ref{ch:selection}) --- all fitted on
training dates, all frozen across the boundary.

\section{Embargoes: quarantine at the border}
\label{sec:embargo}

Suppose training labels at time $t$ involve data through $t + H$ (our
synthetic auxiliary targets reach $H = 40$ steps forward; the target itself
reaches 20). Then a training row within $H$ steps of the validation
boundary \emph{overlaps} the validation period: the model has literally
been fit on statistics of rows it will be evaluated on. The cure is an
\textbf{embargo}: discard a buffer of width $\ge H$ between the last
training timestamp and the first validation timestamp.

Sizing rule: the embargo must cover \emph{the longest forward reach of
anything in the pipeline} --- label construction, forward-shifted auxiliary
targets, rolling labels --- plus the longest \emph{backward} reach of
validation-period features into training rows if those features are shared.
In our setup one full trading day (968 steps) safely dominates the 40-step
maximum reach, so a one-day embargo suffices and costs almost nothing. In
the DRW crypto competition, the winner's cross-validation used purged
group splits with a one-group gap for the same reason
(Chapter~\ref{ch:validation}).

\begin{warnbox}[: the leak may be in the \emph{labels}, not the features]
Teams audit features obsessively and then compute a synthetic training
target with \code{shift(-20)} that crosses the fold boundary. Nothing in
the \emph{feature} pipeline is wrong; the \emph{label} pipeline leaks.
Audit both directions: what do inputs reach backward into, and what do
targets reach forward into?
\end{warnbox}

\section{The audit protocol}

Static review catches construction bugs; only \emph{behavioral} tests catch
the leaks you did not think of. Three tests, in increasing order of cost:

\subsection{The shift test}

Recompute the suspect feature (or the whole feature matrix) shifted forward
by one timestep, retrain quickly, re-score. \textbf{Genuine predictive
signal degrades gracefully} under a one-step delay --- the world is
autocorrelated. \textbf{Leaked signal collapses}, because its value lived
in exact contemporaneity with the future. A validation score that halves
under a one-step shift is a red alert.

\subsection{The shuffle test}

Shuffle the within-day order of rows (keeping (feature, label) pairs
intact), retrain, re-score. A sequence model whose score \emph{survives}
within-day shuffling is not using temporal structure --- fine if that is
expected, alarming if you believed its edge was temporal. A score that
\emph{improves} under shuffling usually indicates the day-boundary handling
was broken all along.

\subsection{The too-good test}

Maintain a standing ceiling estimate (Chapter~\ref{ch:problem}). Any result
above it triggers a leak hunt \emph{before} celebration. Our operating
threshold in this competition: validation $\Rw > 0.03$ meant ``find the
bug.'' It fired twice; both times it was right.

\begin{jsbox}[: two representative catches]
(1) A rolling per-symbol statistic briefly computed with a centered window
during a refactor --- one keyword argument --- produced $+0.02$ of
fictitious score. Caught by the too-good threshold, confirmed by the shift
test (score collapsed), fixed by the trailing-only helper API.
(2) An early draft of the synthetic auxiliary targets crossed the
train/validation boundary through their $\code{shift}(-40)$ reach ---
invisible in feature audits since no \emph{feature} was wrong. Caught in
label-direction review; fixed with the one-day embargo, which we then made
a permanent default of the dataset precomputation script.
\end{jsbox}

\section{Leakage in the large: experiment hygiene}

Information also leaks between \emph{experiments}, through you. Each time
you evaluate a candidate on the validation tail and keep what scored
higher, you fit the tail a little. Over a month of experiments this
selection pressure fabricates $10^{-3}$-scale improvements out of nothing
--- at exactly the effect size you care about. Mitigations, all cheap:

\begin{itemize}
\item \textbf{Two-tier validation}: iterate on window $A$; only finalists
  get scored on window $B$; window $B$ verdicts are final and rare.
\item \textbf{Pre-registration}: one line in the experiment log ---
  \emph{hypothesis, metric, decision rule} --- written before the run.
  This kills post-hoc rationalization (``it lost overall but look at the
  last 20 days'') at negligible cost.
\item \textbf{Seed-replication requirements} for any claimed effect under
  $0.001$ (Chapter~\ref{ch:problem}).
\item \textbf{An honest graveyard}: our \code{WHAT\_DIDNT\_WORK.md}
  records every discarded direction with its evidence. Its quiet function
  is anti-leak: it stops re-litigating ideas whose earlier ``promising''
  runs were, in hindsight, selection noise.
\end{itemize}

\section{Checklist}

Before any result leaves your desk:

\begin{enumerate}
\item Every feature typed trailing / contemporaneous / forward, with the
  forward set empty of model inputs.
\item Lags implemented as date-shifted joins; rolling ops trailing by
  construction.
\item All preprocessing statistics fit on training dates only; frozen
  across the boundary; online updates use only predicted-past days.
\item Embargo at least as wide as the pipeline's longest forward reach,
  labels included.
\item Shift test passed (graceful degradation), shuffle test consistent
  with the model's supposed mechanism, score below the standing ceiling.
\item Experiment pre-registered; effects under $0.001$ replicated across
  seeds and a second window.
\end{enumerate}

Six lines that cost an afternoon to set up bought us months of trustworthy
measurement. Nothing else in this book works without them: every number in
every later chapter is credible exactly because it survived this one.
